%-------------------------------------------------------------------------------
%	SECTION TITLE
%-------------------------------------------------------------------------------
\cvsection{Projects}


%-------------------------------------------------------------------------------
%	CONTENT
%-------------------------------------------------------------------------------
\begin{cventries}

%---------------------------------------------------------
  \cventry
    {Selected personal ML \& LLM projects, open source at github.com/josh-gree}
    {}
    {} % Location
    {2023 -- 2026} % Date(s)
    {
      \begin{cvitems}
        \item {\textbf{probe-lab}: toolkit for LLM activation probing. Extracts per-layer residual-stream activations with PyTorch and Hugging Face Transformers and fits linear probes against prompt metadata for interpretability experiments.}
        \item {\textbf{weight-corruption}: experimental framework measuring how weight perturbations (noise, zeroing, shuffling, quantisation) degrade language modelling and behaviour across the Pythia model family, with GPU evaluation on Modal.}
        \item {\textbf{birdclef-2026} and \textbf{wm}: bird-audio classification in PyTorch (timm, ONNX embeddings) trained on Modal GPUs with Weights \& Biases tracking; wrote \texttt{wm}, a small framework that snapshots every experiment run as a PR for reproducibility.}
        \item {\textbf{news-knowledge-graph}: LLM extraction pipeline turning news articles into a knowledge graph using DSPy programs and GLiNER2 for entity extraction, temporal tagging and Wikidata entity linking, orchestrated with Prefect.}
        \item {\textbf{mobile-agents}: GitHub Actions workflow built on the Claude Agent SDK that plans and implements issues autonomously and opens pull requests. Earlier: RAG-style semantic search over news with DistilBERT embeddings in Postgres/pgvector.}
      \end{cvitems}
    }

%---------------------------------------------------------
\end{cventries}
